\section{Approximate modes can certify the measured curve}\label{sec:curve-certificate}
For the positive difference measure $d\mu$ on $[0,1]$, write
\[
 e_j=\int\lambda^j d\mu,\qquad
 Y_0-Y_j=\int q_j(\lambda)d\mu,\quad
 q_j=\sum_{k=0}^{j-1}\lambda^k.
\]
Suppose node neighborhoods of radius at most $\epsilon$ are assigned
disjointly, with union $S$, and the polynomial certificate gives
$\mu(S^c)\le\eta/d^2$. Collapse each neighborhood's true mass to its
center, omit the exterior mass, and anchor at the same $Y_0$. On $[0,1]$,
$0\le q_j\le j$ and $|q_j'|\le j(j-1)/2$. Splitting the integral gives
\[
 |Y_j-\widetilde Y_j|\le
 \frac{j(j-1)}2\epsilon\mu(S)+\frac{j\eta}{d^2}.
\]
This finite-grid curve bound needs no cutoff away from $\lambda=1$.
Estimated neighborhood weights with total absolute error $\gamma$ and
anchor error $\epsilon_0$ add $j\gamma+\epsilon_0$. True neighborhood
masses cannot silently be identified with fitted weights. An atom at 1 is
an anchored linear finite-window surrogate, not a finite-amplitude plateau.
Plateau inference still pays the singular factor $(1-\lambda)^{-1}$.

For continuous grid coordinate $t\ge 1$, use
$q_t=(1-\lambda^t)/(1-\lambda)$ with $q_t(1)=t$.
Its maximum is $t$ and a Lipschitz bound is
$\max\{1,t(t-1)/2\}$. Indeed
\[
 q_t=t\int_0^1[\lambda+(1-\lambda)u]^{t-1}\,du;
\]
differentiation bounds the maximum derivative by its endpoint value 1
for $1\le t\le 2$, and by $t(t-1)/2$ for $t\ge 2$.
The same split proves a continuous-window bound on this range. A window
touching 0 with unrestricted fast modes requires an additional modulus.

\section{Simultaneous noise coverage permits selected certificates}
Let $\widehat Y=Y+\Sigma^{1/2}g$, where $g$ is an eight-dimensional
standard Gaussian and the full covariance $\Sigma$ is known. On the
single event $\|g\|^2\le 24$,
\[
 |a^{\rm T}(\widehat Y-Y)|\le\sqrt{24a^{\rm T}\Sigma a}
 \quad\hbox{for every }a.
\]
The Gaussian moment generating function gives
$\Pr\{\|g\|^2\ge n+2\sqrt{nx}+2x\}\le e^{-x}$.
For $n=8,x=3$ its threshold is less than 24, so the event has coverage
at least $1-e^{-3}>0.95$. It holds for vectors selected from the same data.
In particular, if $c$ selects a polynomial after observing the data,
$\eta=c^{\rm T}Hc=a(c)^{\rm T}Y$ has the simultaneous upper bound
\[
 \eta\le\widehat\eta+\sqrt{24a(c)^{\rm T}\Sigma a(c)}.
\]
Estimated covariance, heavy tails and unmodeled drift need a separately
calibrated confidence event. No Gaussian coverage is assigned to them.

Differences reuse observations: $C=D\Sigma D^{\rm T}$.
For the declared fixture $\Sigma_{ij}=\sigma^2 2^{-|i-j|}$,
$\sigma=10^{-6}$, each difference variance is $\sigma^2$.
The same event therefore bounds every moment entry error by $5\sigma$.
For a size-$r$ Hankel matrix,
\[
 \widehat H_r-5r\sigma I\succ 0
 \quad\Longrightarrow\quad H_r\succ 0.
\]
This certifies at least $r$ visible nodes for every compatible positive
measure. Failure is unresolved, not proof of smaller exact rank.
Noise cannot exclude an arbitrarily weak extra positive mode.

The exact two-mode fixture has $\alpha=4/5$, $B=(1/10,1/10)$,
$\lambda=(1/4,3/4)$ and eight point samples. It passes the size-two lower
certificate. A seeded covariance-aware calculation uses 2000 catalogs in
each of five declared cases: separated, nearly coincident, weak second,
very slow and constant. The separated case certifies a minimum of two;
the others remain unresolved at two. These are new inference checks,
separate from the historical Rev 6 benchmark or observational data.

\subsection{Positive support constraints inside the confidence event}
For eight observations let $e_j=Y_j-Y_{j+1}$, $0\le j\le 6$.
The exact degree-six positive difference-measure constraints on $[0,1]$ are
\[
 \begin{gathered}
 H_4(e)=(e_{i+j})_{i,j=0}^3\succeq 0,\\
 L_3(e)=(e_{i+j+1}-e_{i+j+2})_{i,j=0}^2\succeq 0.
 \end{gathered}
\]
Intersect these with
$\|\Sigma^{-1/2}(Y-\widehat Y)\|_2\le\sqrt{24}$.
These conditions are sufficient as well as necessary by the one-variable
representation of a nonnegative degree-six polynomial as sums of squares
plus $\lambda(1-\lambda)$ times sums of squares.
Maximizing a selected polynomial's linear moment functional over this
convex set gives a simultaneous upper certificate; the complete conic
optimization is a specified method, not an executed optimum here.

A finite implemented dual method supplies usable bounds. For positive
Gram multipliers $P,Q$, write $b=a+H^*P+L^*Q$. Then every feasible $Y$ obeys
\[
 a^TY\le b^T\widehat Y+\sqrt{24b^T\Sigma b}.
\]
The minus signs give the corresponding lower certificate. A catalogue of 40
exact Gram choices, including zero, is evaluated with outward rational
square roots. Selection within this catalogue retains the same confidence
event. An explicit constant-curve fixture improves the unconstrained
ellipsoid's positive upper bound on $-e_0$ to exactly 0, with a feasible
true curve; a separate drift fixture has an exact infeasibility witness.

If the class also requires positive amplitudes and a nonnegative plateau,
use the stronger positive-$Y$ support constraints
\[
 (Y_{i+j+1})_{i,j=0}^3\succeq 0,\qquad
 (Y_{i+j}-Y_{i+j+1})_{i,j=0}^3\succeq 0.
\]
An atom of the difference measure at $\lambda=1$ is a valid anchored
finite-window linear surrogate, not a finite-amplitude plateau model.
For a fitted curve, two linear optimizations per grid point bound
$\pm(Y_j-\widetilde Y_j)$ directly over the chosen intersection; this avoids
identifying true neighborhood masses with fitted weights. The supplement
gives the full formulations, exact feasibility witnesses and finite dual
implementation separately.


\subsection{A window mismatch is a model error}
Point-sample Hankel formulas also apply to common translated bin windows:
each exponential acquires a common positive amplitude factor and retains
its node. Unequal windows generally do not. For one true node $1/2$ and amplitude $B=1/5$,
an alternating positive two-point window factor 1 and 5/4 gives the
four-point difference determinant $-27/6400$, whereas a common window
gives 0. This is a failure of the point-sample model, not evidence for a
positive two-mode model. Appendix~\ref{app:statistics}'s window correction
is required before the main-text single-mode and rank tests are applied.

\section{Coarse winding and fine final readout use different ensembles}
The earlier twenty-fine-mean design is sufficient but spends final
precision on every winding prefix. Assume the controlled initial phase
is known. Use three intermediate prefixes per history with two coordinates
each:12 coarse means with coordinate tolerance $d_c=.05$ and 8100 fresh
preparations per mean. Use only the two final coordinates of each history
at $d_f=5\,10^{-11}$: four fine means, each with
$5.432\,10^{21}$ fresh preparations. Hoeffding and the union bound give
\[
 \begin{aligned}
 P_c&\le 24e^{-(8100)d_c^2/2}<.000961568,\\
 P_f&\le 8e^{-(5.432\,10^{21})d_f^2/2}<.008999751.
 \end{aligned}
\]
The total statistical failure is below $.009961318<.01$ and the sufficient
preparation count is
\[
 21\,728\,000\,000\,000\,000\,097\,200
 \simeq 2.1728\,10^{22}.
\]
The inherited amplitude lower bound, controlled prefix spacing and
separate coarse systematic allowance keep every unwrapping increment
strictly below $\pi$. Fine final sampling retains the earlier statistical
phase allowance below 0.2 nanoradians. Unknown initial phase requires
additional sampling. Every original preparation, leakage, readout,
timing, reference, energy and model requirement remains in force.
This failure probability is conditional on the stated systematic bounds;
random calibration failures must receive their own probability allocation.
These are sufficient mathematical resource estimates, not achieved
apparatus performance or experimental lower bounds.
